\documentclass[10pt,a4paper]{amsart}
\usepackage[margin=19mm]{geometry}
\usepackage{amsmath,amssymb,mathtools}
\usepackage[scr=rsfs,cal=euler]{mathalfa}
\usepackage{xfrac}
\usepackage[unicode,hidelinks,bookmarks=false]{hyperref}
\newcommand{\ie}{i.e.\ }
\newcommand{\cf}{cf.\ }
\newcommand{\resp}{respectively\ }
\newcommand{\Subsection}{\subsection}
\providecommand{\nobreakdash}{\nobreak\hbox{-}}
\setlength{\emergencystretch}{2em}
\multlinegap0pt
\usepackage{amssymb}
\usepackage{enumerate}
\usepackage{mathtools}
\newtheorem{lem}{Lemma}
\providecommand{\Linf}{\mathcal{L}_{\infty}}
\providecommand{\Lneg}{\mathcal{L}_{\infty}^-}
\providecommand{\Lpos}{\mathcal{L}_{\infty}^+}
\providecommand{\Ns}{N \times I}
\providecommand{\Rn}{\mathbb{R}^n}
\providecommand{\Rnp}{D_{\Omega}}
\title{On an evolutionary equation involving the nonlocal infinity Laplacian}
\author{Frida Fejne}
\date{Source: arXiv:2609.11311v1 (CC BY 4.0)}
\begin{document}
\maketitle

\noindent\emph{Source and licence.} On an evolutionary equation involving the nonlocal infinity Laplacian. Version arXiv:2609.11311v1 (CC BY 4.0), distributed by arXiv under the Creative Commons Attribution 4.0 International licence, http://creativecommons.org/licenses/by/4.0/, as recorded in the arXiv licence metadata for this exact version and shown on the abstract page https://arxiv.org/abs/2609.11311v1. Full licence notice: \texttt{licenses/arxiv-2609.11311-CC-BY-4.0.txt}.\par\smallskip

\noindent\textit{Editorial scope.} Counted unit: the article's lem lem:supsubissub (source0.tex lines 697--703). The article's complete original proof of it is delivered with it (source0.tex lines 704--764, one \texttt{proof} environment of 61 lines). No mathematics is written, repaired or re-derived: the statement and the proof are the article's own lines, in the article's own notation, and no retained line is substituted or re-flowed.  Retained uncounted as the source-local context the proof needs: lines 89--91; lines 186--188; lines 190--192; lines 243--248; lines 320--326; lines 1332--1342.  Nothing was omitted from any retained span, so the package carries no omission record.  No retained line contains a citation, so the wrapper carries no bibliography and no bibliography entry.  No retained line contains a figure environment, an image-inclusion command or drawing code, so no image is delivered and the package declares no asset dependency.  Editorial formatting only, in addition: an AMS \texttt{amsart} preamble that restates the article's own declarations and macros used by the retained lines, namely the environments \texttt{lem} on separate counters, together with the standard packages those lines need.  The retained environments print in this document as Lemma 1 in order of appearance, whereas the article prints the same items with the numbering of its own full document.  The complete native source of the article as distributed in the arXiv e-print for this exact version is bundled unchanged as \texttt{sources/210006/source0.tex}.
\begin{align} \label{eq:Luf} 
 \frac{\partial u}{\partial t} = (\mathcal{L}_{\infty})^{\alpha}u   \textup{ in }  \Omega,
\end{align}

\begin{equation}\label{eq:u_upp}
    u^*(x,t) := \lim_{\epsilon \to 0} \sup_{(y,s) \in B_{\epsilon}(x,t)} u(y,s),
\end{equation}

\begin{equation}\label{eq:u_low}
    v_*(x,t) := \lim_{\epsilon \to 0} \inf_{(y,s) \in B_{\epsilon}(x,t)} v(y,s).
\end{equation}

\begin{align}\label{eq:varphiN}
     \varphi_{N}(x,t)= \begin{cases}
                      \varphi & \textup{if  } \ (x,t) \in \Ns,\\
                      u & \textup{if  } (x,t) \in \ \Rnp   \setminus \Ns.
            \end{cases}
 \end{align} 

\begin{lem}\label{lem:Linfusc}
    Let $0<\alpha \leq 1$, $u$ be upper semicontinuous, $\Ns \subset \Omega$ and $\varphi \in C^1(\Ns)$. Then $\Linf \varphi_{N}$ is upper semicontinuous in $\Ns$, that is for $(x_0,t_0) \in \Ns$ we have that
    \begin{equation*}
        \limsup_{(x_k,t_k) \to (x_0,t_0)} \Linf \varphi_{N}(x_k,t_k) \leq \Linf \varphi_{N}(x_0,t_0),
    \end{equation*}
    where $\varphi_{N}$ is defined as in \eqref{eq:varphiN}. If $v$ is lower semicontinuous, then $\Linf \varphi_{N}$ is lower semicontinuous in $\Ns$.
\end{lem}

\begin{lem}\label{lem:strictmax}
    Let $\mathcal{A}$ be a family of upper semicontinuous functions. We let 
    \begin{equation*}
        u(x,t):= \sup_{v \in \mathcal{A}} v(x,t),
    \end{equation*}
and define $u^*$ as in \eqref{eq:u_upp}. We consider a point $(x_0, t_0) \in \Omega.$ Let $\Ns $ be a neighborhood of $(x_0,t_0)$ such that $(x_0,t_0)\in \Ns \subset \Omega$. Furthermore, let $\varphi \in C^1(\Ns)$ be a function that touches $u^*$ strictly from above at $(x_0,t_0)$. Then there exists a neighborhood $Q_r= B_r(x_0,t_0) \times (t_0-r,t_0+r)$ of $(x_0,t_0)$ such that $Q_r \subset \subset \Ns$, and a sequence of points $(x_k,t_k) \in Q_r$, such that $(x_k,t_k)\to (x_0, t_0)$, and functions $u_k \in \mathcal{A}$ such that 
    \begin{equation}\label{eq:supoverN}
         \sup_{Q_r} u_k-\varphi = u_k(x_k,t_k)-\varphi(x_k,t_k),
    \end{equation}
     and $u_k(x_k,t_k) \to u^*(x_0, t_0)$.
\end{lem}

\begin{lem}\label{lem:supsubissub}
    Let $\mathcal{A}$ be a family of viscosity subsolutions (resp. viscosity supersolutions) of \eqref{eq:Luf} and define 
    \begin{equation}\label{eq:sup_u}
        u(x,t):= \sup_{w \in \mathcal{A}} w(x,t) \quad (\textup{resp. } v(x,t):= \inf_{w \in \mathcal{A}} w(x,t) ).
    \end{equation}
Then $u^*$ (resp $v_*$) is a viscosity subsolution (resp. supersolution) of $\eqref{eq:Luf}$ in $\Omega$, where $u^*$ and $v_*$ are defined in \eqref{eq:u_upp} and \eqref{eq:u_low}, respectively.
\end{lem}

\begin{proof}
    We provide the proof for subsolutions. We consider $(x_0, t_0) \in \Omega$, a neighborhood $\Ns$ of $(x_0,t_0)$ such that $ \Ns \subset \Omega$ and let we let $\varphi \in C^1(\Ns)$ be a test function that touches $u^*$ strictly from above at $(x_0,t_0)$. From Lemma \ref{lem:strictmax} we know that that there exists a cylinder $Q_r=B_r(x_0,t_0) \times (t_0-r,t_0+r)$, such that $Q_r \subset \subset \Ns$, and points $(x_k,t_k) \in Q_r$, such that $(x_k,t_k)\to (x_0, t_0)$, and functions $u_k \in \mathcal{A}$ such that 
  \begin{equation}\label{eq:supoutsideuk}
         \sup_{Q_r } (u_k-\varphi) = u_k(x_k,t_k)-\varphi(x_k,t_k),
    \end{equation}
     and $u_k(x_k,t_k) \to u^*(x_0, t_0)$. We note that $u_k(x_k,t_k)\leq u^*(x_k,t_k) < \varphi(x_k,t_k)$ for each $(x_k,t_k) \neq (x_0,t_0)$. However, from \eqref{eq:supoutsideuk} it follows that the function $\varphi-M_k$ touches $u_k$ from above at $(x_k,t_k)$, where
     $$
     M_k = \inf_{Q_r} (\varphi-u_k) = -\sup_{Q_r } (u_k-\varphi) = \varphi(x_k,t_k)-u_k(x_k,t_k).
     $$
     We next define $\varphi^k=\varphi-M_k$ and
    \begin{align*}
     \varphi_{r }^k= \begin{cases}
                      \varphi-M_k & \textup{if  } \ (x,t) \in Q_r ,\\
                      u_k & \textup{if  } (x,t) \in \ \Rnp \setminus Q_r,
            \end{cases}
 \end{align*}
and recall that 
\begin{align*}
     \varphi_{r }= \begin{cases}
                      \varphi & \textup{if  } \ (x,t) \in Q_r ,\\
                      u^* & \textup{if  } (x,t) \in \ \Rnp \setminus Q_r.
            \end{cases}
 \end{align*}
We introduce the notation
    \begin{equation*}
        L^+_{r}(x,t) = \sup_{y \in B_r(x_0)} \frac{\varphi(y,t)-\varphi(x,t)}{|y-x|^{\alpha}} \quad \textup{and } \quad L^-_{r}(x,t) = \inf_{y \in B_r(x_0)} \frac{\varphi(y,t)-\varphi(x,t)}{|y-x|^{\alpha}},
    \end{equation*}
and
    \begin{equation*}
        L^+_{r^c}(x,t) = \sup_{y \in \Rn \setminus B_r(x_0)} \frac{u_k(y,t)-\varphi(x,t)+M_k}{|y-x|^{\alpha}} \quad \textup{and } \quad L^-_{r^c}(x,t) = \inf_{y \in \Rn \setminus B_r(x_0)} \frac{u_k(y,t)-\varphi(x,t) +M_k}{|y-x|^{\alpha}},
    \end{equation*}
    and we note that for $(x,t) \in Q_{r/2}$ we have that
\begin{align*}
     L^+_{r^c}(x,t) &\leq \sup_{y \in \Rn \setminus B_r(x_0)} \frac{u_k(y,t)-\varphi(x,t)}{|y-x|^{\alpha}} + \sup_{y \in \Rn \setminus B_r(x_0)} \frac{M_k}{|y-x|^{\alpha}} \\
     &= \sup_{y \in \Rn \setminus B_r(x_0)} \frac{u_k(y,t)-\varphi(x,t)}{|y-x|^{\alpha}} + \frac{M_k}{\inf_{y \in \Rn \setminus B_r(x_0)}|y-x|^{\alpha}} \\
     & \leq \sup_{y \in \Rn \setminus B_r(x_0)} \frac{u^*(y,t)-\varphi(x,t)}{|y-x|^{\alpha}} + \frac{2^{\alpha} M_k}{r^{\alpha}}.
\end{align*}
Furthermore,
\begin{align*}
    L^-_{r^c}(x,t) &\leq \inf_{y \in \Rn \setminus B_r(x_0)} \frac{u_k(y,t)-\varphi(x,t)}{|y-x|^{\alpha}} + \sup_{y \in \Rn \setminus B_r(x_0)} \frac{M_k}{|y-x|^{\alpha}} \\
     &= \inf_{y \in \Rn \setminus B_r(x_0)} \frac{u_k(y,t)-\varphi(x,t)}{|y-x|^{\alpha}} + \frac{M_k}{\inf_{y \in \Rn \setminus B_r(x_0)}|y-x|^{\alpha}} \\
     & \leq \inf_{y \in \Rn \setminus B_r(x_0)} \frac{u^*(y,t)-\varphi(x,t)}{|y-x|^{\alpha}} + \frac{2^{\alpha} M_k}{r^{\alpha}}.
\end{align*}
For $(x,t) \in Q_r$ we have that
    \begin{align*}
        \Lpos \varphi_{r}^k(x,t) &= \max(L^+_{r^c}(x,t), L^+_{r}(x,t)), \\
        \Lneg \varphi_{r}^k(x,t) &= \min(L^-_{r^c}(x,t), L^-_{r}(x,t)).
    \end{align*}
Thus, for $k$ large enough such that $(x_k,t_k) \in Q_{r/2}$ it follows that
\begin{equation*}
    \Linf \varphi^k_{r}(x_k,t_k) \leq \Linf \varphi_{r}(x_k,t_k)+ \frac{2^{\alpha} M_k}{r^{\alpha}},
\end{equation*}
where the last term goes to zero $k \to \infty$ since $M_k$ goes to zero as $k \to \infty$. Since each $u_k$ is a subsolution in the viscosity sense it follows that
    \begin{align*}
       0 &\geq \lim_{k \to \infty} \left(\frac{\partial \varphi^k}{\partial t}(x_k, t_k) - \Linf \varphi_{r }^k(x_k, t_k) \right)\\
       &\geq \lim_{k \to \infty} \left(\frac{\partial \varphi}{\partial t}(x_k, t_k) - \Linf \varphi_{r }(x_k, t_k) -\frac{2^{\alpha} M_k}{r^{\alpha}}\right)\\
       &\geq \frac{\partial \varphi}{\partial t}(x_0, t_0) - \Linf \varphi_{r}(x_0, t_0)\\
       & \geq  \frac{\partial \varphi}{\partial t}(x_0, t_0) - \Linf \varphi_{N }(x_0, t_0),
    \end{align*}
where we used Lemma \ref{lem:Linfusc} in the third inequality. The last inequality follows from the fact that $\Linf \varphi_{r }(x_0, t_0) \leq  \Linf \varphi_{N }(x_0, t_0)$ since $\varphi \geq u^*$ in $\Ns$. Thus, we have proved that $u^*$ is a viscosity subsolution.
\end{proof}

\end{document}
